\begin{abstractbox}
\textbf{\color{iopBlue}Abstract.} We report a systematic study of vortex pinning enhancement in niobium nitride (NbN) superconducting thin films following irradiation with 220~MeV Au$^{7+}$ heavy ions. Four film sets, grown by reactive DC magnetron sputtering to a thickness of 25~nm on c-plane sapphire, were irradiated at fluences spanning zero to $1\times10^{12}$~ions~cm$^{-2}$ and characterised by resistive transport, field-dependent critical current measurements, and magneto-optical flux imaging. The critical temperature decreases monotonically with fluence, from $15.8$~K in the pristine film to $10.7$~K at the highest dose, consistent with pair-breaking by irradiation-induced disorder. The zero-field critical current density at $4.2$~K, however, is non-monotonic: it rises from $6.8\times10^{6}$~A~cm$^{-2}$ in the pristine film to a maximum of $12.1\times10^{6}$~A~cm$^{-2}$ at an intermediate fluence of $5\times10^{11}$~ions~cm$^{-2}$ before falling to $7.9\times10^{6}$~A~cm$^{-2}$ at the highest dose, where superfluid-density suppression outweighs the added pinning capacity. Fits of the field-dependent pinning force density to the Dew-Hughes scaling form yield exponents $(p,q) = (0.51 \pm 0.04,\, 1.98\pm0.07)$ at the optimal dose, consistent with a crossover from weak point-pinning to a regime dominated by extended, surface-normal damage columns. Magneto-optical imaging corroborates this picture, showing a marked suppression of dendritic flux avalanches in the irradiated films relative to the pristine reference. These results identify an irradiation window that nearly doubles the critical current density of NbN films without unacceptable suppression of $T_c$, of direct relevance to superconducting nanowire single-photon detectors and superconducting radio-frequency cavities that require both high pinning strength and a robust transition temperature.
\end{abstractbox}
